\begin{exercise}{132}
Suppose that you are given a cocartesian fibration $p\colon\mathcal E\to\mathcal C$. Recall that we
have constructed for every edge $f\colon\Delta^1\to\mathcal C$ a functor $\mathcal E_x\times\Delta^1\to\mathcal E$, whose restriction to
$\mathcal E_x\times\{1\}$ is given by $f_!$. Show that this functor may equivalently be constructed by considering
the diagram
\[
\begin{tikzcd}[column sep=large,row sep=large]
 & \mathcal E_x\times\Delta^{\{0\}}\arrow[r]\arrow[d] & \mathcal E\arrow[d] \\
 \mathcal E_x\times\{1\}\arrow[r] &
 \mathcal E_x\times\Delta^1\arrow[r,"f\circ\mathrm{pr}"]\arrow[ur,dashed] & \mathcal C
\end{tikzcd}
\]
and show the existence of a dashed arrow having all the above properties.

\begin{proof}
By Proposition~3.2.13, postcomposition gives a cocartesian fibration
\[
 p_*\colon\mathrm{Fun}(\mathcal E_x,\mathcal E)
 \longrightarrow\mathrm{Fun}(\mathcal E_x,\mathcal C).
\]
The constant family of edges $f$ is an edge in the target from the
constant $x$-functor to the constant $y$-functor, where $y=f(1)$.
Lift it to a $p_*$-cocartesian edge starting at the fibre inclusion
$\mathcal E_x\to\mathcal E$.
Its adjoint $H\colon\mathcal E_x\times\Delta^1\to\mathcal E$ fills the
displayed square, restricts at $0$ to that inclusion, and lies over
$f\circ\mathrm{pr}$.
The pointwise characterization in Proposition~3.2.13 makes each
$H(z,-)$ a $p$-cocartesian lift of $f$.
Its restriction at $1$ therefore lands in $\mathcal E_y$.

Equivalently, adjunction gives a section
$\mathcal E_x\to\mathrm{Fun}^{\mathrm{cc}}_f(\Delta^1,\mathcal E)$
of evaluation at $0$.
Evaluation at $1$ is consequently exactly the functor $f_!$ constructed
in Proposition~3.3.1 using this section.
\end{proof}
\end{exercise}
